\section{Part II: The Financial Statements and Economic Reality}

\subsection{18. Accounting is a model of events, not a photograph of cash}


Financial statements record economic claims under specified measurement, recognition and presentation rules. They are indispensable, but the reader must understand what they represent. Revenue can be earned before collection; an expense can be recognized before payment; a capital asset can consume cash this year while contributing depreciation expense across future periods. Estimates and judgments enter expected credit losses, useful lives, contract revenue, impairments and contingent obligations.

The founder needs 2 ledgers: \textbf{statutory financial statements} prepared under applicable standards and \textbf{managerial economic analysis} designed to answer operating questions. The latter may value founder labor at a market replacement cost, split cohorts differently or forecast future committed delivery. Do not quietly modify statutory numbers to fit a management narrative. Reconcile adjusted figures to their reported starting points.

\subsection{19. The balance sheet: assets are financed by claims}


At every reporting date: $\text{Assets} = \text{Liabilities} + \text{Equity}$. Assets may include cash, receivables, inventory, equipment, prepaid expenses, investments and validly recognized intangible assets. Liabilities may include suppliers, tax, payroll, debt, leases, customer advances and provisions. Equity is the residual accounting claim, not a reliable estimate of business market value.

A founder can create a valuable software business with little book equity because much internally generated know-how is expensed. Another may own expensive property that produces poor returns. Book value is not automatically liquidation value and neither equals a negotiated share price. Ask what each balance-sheet item can actually \textbf{earn, convert to cash or require in settlement}.

\subsection{20. The income statement: an account of period performance}


An income statement usually progresses from recognized revenue through direct and operating costs, depreciation and amortization, interest, tax and profit. The exact headings depend on jurisdiction and reporting standards. Gross margin can reveal direct delivery economics, but allocation rules matter. A consulting company calling all engineers 'administration' may report inflated gross margin. A software company expensing engineers but capitalizing selected development costs may show different operating profit from a peer with different legitimate policy choices.

The founder should maintain operational versions of contribution by program or customer after \textbf{incremental delivery, support, refunds, payment fees and acquisition}. Net profit then reconciles company-level staff, rent, depreciation, finance and taxes. A profitable unit is not yet proof of a profitable company; an accounting loss is not necessarily evidence that every customer destroys value.

\subsection{21. The cash-flow statement: where money was generated and consumed}


The cash-flow statement groups cash into operations, investing and financing under the applicable accounting standards. Operating cash reflects receipts and payments related to the business, often presented indirectly by adjusting profit for noncash items and working-capital changes. Investing cash generally includes equipment purchases, investments and disposals. Financing cash generally includes borrowing and repayment, equity issuance and distributions, with presentation differences for some interest/dividend items across frameworks.

\textbf{Cash from operations is not free cash flow.} A company can generate ₹20 lakh operating cash, spend ₹35 lakh maintaining equipment, and therefore need financing. Conversely, operating cash can spike because customers prepay services that still must be delivered. Separate \textit{sustainable recurring cash generation} from 1-time timing movements.

\subsection{22. A complete 3-statement transaction}


Assume an unlevered educational microcompany starts with \textbf{₹5,00,000 founder equity}, then collects \textbf{₹3,00,000} for teaching to be provided over 3 months. At collection: cash increases by ₹3,00,000 and a contract liability / deferred revenue increases by ₹3,00,000. No tuition revenue has yet been earned. Suppose the 1st month delivers 1/3 of the promised service and spends ₹70,000 on faculty and related monthly expenses, paid in cash. At month-end the company recognizes ₹1,00,000 revenue and ₹70,000 expense, yielding ₹30,000 profit.

The balance sheet is: \textbf{cash ₹7,30,000; liabilities ₹2,00,000; equity ₹5,30,000}. The identity balances. Operating cash flow for the month is ₹2,30,000: ₹3,00,000 customer receipt less ₹70,000 cash expense. Under an indirect bridge, ₹30,000 profit plus the ₹2,00,000 increase in deferred revenue also equals ₹2,30,000. Founder contributed ₹5,00,000 financing cash earlier. The bank balance is not ₹7,30,000 of available profit: ₹2,00,000 of future service remains owed, and further direct costs will arise. Tax and other obligations are omitted solely for clarity.

\subsection{23. Revenue is not cash, bookings, invoices or billings}


A '₹30 lakh annual contract' can mean signed bookings, a legal right to bill, an issued invoice, a collection, or revenue recognized from fulfilled performance. Their timings and cancellation conditions differ. For subscriptions, annually prepaid receipts are often deferred and recognized across the service period. For projects, revenue may be recognized over time or at a point in time depending on contract terms and applicable standards. \textbf{IFRS 15} uses a performance-obligation framework; firms applying Indian AS, Ind AS or other rules must apply the appropriate regime rather than assume all IFRS details bind them.

An increase in deferred revenue can temporarily improve cash while indicating future obligations. Revenue recognized but uncollected increases receivables and may reduce liquidity. In board reporting, show \textbf{bookings, billings, recognized revenue, cash collections, and remaining performance obligations} separately if they matter to the model.

\subsection{24. Cash and accrual reconciliation: the 4-question audit}


For each material difference between reported profit and cash, ask: \textbf{Was the transaction paid? Was it earned or consumed? Was it capitalized? Does it belong to another period?} An unpaid customer invoice creates receivables. A prepaid software subscription creates a prepaid asset, expensed as used. An equipment purchase creates an asset and an immediate investing cash outflow, with depreciation later. A share issuance creates financing cash and equity but not revenue.

A founder who cannot explain such differences will misread both growth and solvency. A competent monthly close should produce an understandable bridge from accounting profit to operating cash, from opening to closing cash, and from opening to closing equity. Unexplained 'adjustments' should become questions, not be hidden for presentation purposes.

\subsection{25. Working capital as the financing of the operating cycle}


Operating net working capital can be approximated as \textbf{operating current assets minus non-interest-bearing operating current liabilities}. In an inventory business, purchases consume cash before sale; receivables delay cash after sale; payables may defer settlement. A simple cash conversion cycle is $\mathrm{DIO} + \mathrm{DSO} - \mathrm{DPO}$ (days inventory outstanding, sales outstanding, payables outstanding), subject to meaningful measurement conventions.

Suppose a manufacturer has 45 inventory days, 50 receivable days, and 30 supplier-credit days. Its cycle is 65 days. At an illustrative operating cash cost of ₹20,000 per day, financing 65 days of net operating cost suggests ~₹13 lakh in working-capital exposure-only a rough planning estimate, not a substitute for a full balance-sheet model. DSO, DIO and DPO use different denominators in formal calculations. Growth generally scales this capital need unless terms improve.

\subsection{26. Negative working capital can be an advantage-or a trap}


A healthy subscription business can collect from customers before paying service costs. Some retailers collect immediately and pay suppliers later. This can release capital and improve return on invested capital. But negative working capital is not magic free profit: the business still owes delivery, refunds, customer service or supplier payments. If growth stops, the influx of new prepayments may cease while existing obligations remain.

A school that repeatedly launches new prepaid cohorts and uses new tuition to finish older cohorts is potentially vulnerable if there are no adequate ring-fenced planning reserves and unit economics are weak. Model a \textbf{no-new-sales scenario}: can already-sold commitments be completed from cash and reliable collections? If not, strong bank balances during enrollment periods conceal financial fragility.

\subsection{27. Fixed versus variable costs are decision-specific}


A faculty contractor paid per submitted review may be variable across submissions, but nearly fixed within a committed monthly minimum. Cloud expenditure can include subscriptions, usage and reserved capacity. Marketing spend may be discretionary in the short run but necessary to replace attrition in the long run. Therefore labels should be attached to a \textbf{volume range and time horizon}. A cost is not intrinsically fixed forever.

For decisions, calculate \textbf{avoidable cost}. If a program is cancelled tomorrow, which expenses disappear, which remain due under contract, and what capacity becomes reusable? Allocated headquarters overhead may matter for total profitability but may not change when a single contract is accepted. Conversely, treating every allocated cost as avoidable can cause the company to reject profitable incremental work.

\subsection{28. Contribution, EBITDA, EBIT, NOPAT and free cash flow}


Contribution usually means realized revenue less defined variable or incremental costs; the definition must be explicit. \textbf{EBITDA} is earnings before interest, taxes, depreciation and amortization; it excludes some very real recurring capital consumption and working-capital costs. \textbf{EBIT} incorporates depreciation and amortization. \textbf{NOPAT}, net operating profit after tax, approximates after-tax operating profit independently of financing, often $\mathrm{EBIT} \times (1 - \text{Operating Tax Rate})$. \textbf{FCFF}, free cash flow to the firm, approximates $\mathrm{NOPAT} + \text{Noncash Charges} - \mathrm{Capex} - \Delta \mathrm{NWC}$, with adjustments depending on accounting treatment and business model.

All are useful and none is a complete verdict. Compare a capital-intensive company and a software service cautiously. An equipment-heavy venture can have high EBITDA while spending nearly all of it on maintenance. A growing SaaS business may report modest accounting profit but attractive long-run economics after early product development. A careful founder identifies which expenses will actually recur at steady state.

\subsection{29. Maintenance versus growth capital expenditure}


Maintenance capex preserves the productive ability of existing operations; growth capex adds capacity or functionality. The distinction is economic, not always directly visible in accounts. Replacing worn servers to maintain uptime is maintenance; adding a cluster for new paying demand is growth. A school may need continuous curriculum rewriting simply to remain relevant; this maintenance can appear as operating expense rather than accounting capex.

Management sometimes declares capital expenditure 'growth' to imply higher sustainable free cash flow. Test the claim: \textbf{if no new customers arrive for 3 years, what spend is still necessary to keep current service quality and revenue?} That expenditure is economically maintenance-like regardless of how the accounts present it. A company with low maintenance needs has genuine financial flexibility; 1 that merely omits maintenance from forecasts does not.

\subsection{30. Depreciation and amortization are not fictional costs}


Depreciation allocates an asset's cost over useful periods; amortization does something similar for appropriate intangible assets. Neither is a current cash payment, yet both can approximate real economic consumption when assets wear out or become obsolete. They can diverge substantially from actual replacement economics, especially under technological change and inflation. The founder should forecast both accounting depreciation and the \textbf{cash schedule for eventual replacement}.

A ₹12 lakh machine depreciated over 4 years may show ₹3 lakh annual straight-line depreciation with zero residual value, but actual replacement might be required in year 3 or might never be justified. Do not conflate tax depreciation, accounting depreciation and economic obsolescence. The project investment case should use actual expected cash flows, not depreciation alone.

\subsection{31. Leases can hide capital-like commitments from casual analysis}


Renting an office reduces upfront purchase expenditure but may create an enforceable multi-year payment obligation. Depending on reporting rules, leases can generate recognized right-of-use assets and lease liabilities; \textbf{IFRS 16} addresses this for most lessee leases, subject to exemptions, while a private Indian entity's applicable accounting framework may differ. An entrepreneurial analysis should not confuse 'no property owned' with 'no financing-like commitments'.

Before signing a 5-year campus lease, calculate: refundable deposit; expected escalation; fit-out; obligations at early exit; fixed service costs; tax; restoration; minimum occupancy to cover costs; and stress cash when demand falls. The option to rent rather than buy is valuable only if the rental arrangement actually retains flexibility. A long noncancellable lease can behave like leveraged fixed-capital investment.

\subsection{32. Business versus financial leverage}


\textbf{Operating leverage} arises when fixed operating costs cause profit to change more than proportionately with revenue over a range. \textbf{Financial leverage} arises when borrowing creates fixed interest and principal obligations and magnifies returns or losses to equity holders. A business with both high fixed operating costs and high debt service can be extremely fragile to small demand shocks.

Suppose revenue is ₹100, variable costs ₹40 and fixed operating costs ₹40, so EBIT is ₹20. A 10\% revenue decline with the same variable-cost ratio reduces EBIT to ₹14, a 30\% reduction. Add ₹10 of annual interest and pretax income moves from ₹10 to ₹4, a 60\% decline. These are stylized period figures; tax and cash timing are excluded. \textbf{Leverage magnifies the consequence of being wrong}. It should be attached to stable, fairly predictable cash flows, not assumed to manufacture profitability.

\subsection{33. The debt service coverage ratio is a cash test}


Lenders often examine \textbf{DSCR}, broadly available cash flow for debt service divided by required interest and principal payments, with precise definitions specified by the facility. Suppose a project generates ₹18 lakh cash available for debt service and owes ₹12 lakh in scheduled principal plus interest: $\mathrm{DSCR} = \mathbf{1.50\times}$. If a downside scenario produces ₹10 lakh available, $\mathrm{DSCR} = \mathbf{0.83\times}$, indicating a funding gap absent permitted reserves or remedies.

Use the lender's contractual definition, not an invented EBITDA numerator. Changes in working capital, tax, maintenance capex or lease payments can be consequential. A seemingly comfortable annual DSCR can conceal a quarter with insufficient cash. Model the monthly or quarterly minimum, covenant cure rights and the possibility that an otherwise solvent company loses financing because it breaches a ratio.

\subsection{34. Ratios are useful only when their accounting drivers are visible}


A current ratio above 1 may conceal doubtful receivables; a low current ratio may coexist with reliable recurring customer prepayments and strong collection. A high gross margin may reflect omitted delivery labor. A high ROE may result from debt, not superior operating capability. A low debt-to-equity ratio may conceal leases, guarantees and off-balance contractual commitments.

Instead of declaring 'healthy', walk from the ratio to its source records. \textbf{Who owes the receivable? How old is it? What is the inventory? What is contractually payable? What was recognized rather than collected?} Compare ratios across periods with consistent policy and across genuine peer models. A public listed-company average should not automatically become a covenant for a 10-person founder-led institute.

\subsection{35. Financial statement analysis by stakeholder}


Founders need delivery viability and runway. Equity investors need value creation, dilution and exit possibilities. Lenders prioritize reliable debt-service cash and downside recovery. Employees care about payroll stability and the ability to invest in career development. Buyers care whether the vendor will survive to perform. Donors care about restricted-fund stewardship and mission delivery. Different financial presentations become misleading when 1 stakeholder's favorite measure is offered as a universal proof.

A disciplined financial pack contains \textbf{an income statement, balance sheet, cash-flow statement, 13-week cash forecast, receivables aging, obligations schedule, program/customer economics and relevant covenant or ownership records}. Early firms can implement these in small spreadsheets if bookkeeping is sound. Scaling firms need more reliable close processes and access controls.

\subsection{36. The entrepreneur's quality-of-earnings bridge}


Begin with reported EBITDA or EBIT and reconstruct sustainable earnings. Remove genuine 1-off benefits, reverse unsupported revenue, normalize founder compensation and required staff costs, scrutinize capitalized expenses, include realistic bad-debt and refund allowances, identify related-party price distortions, recognize required maintenance, and separate profitable core activity from temporary subsidies. Adjustments need 2-sided discipline: do not add back every inconvenient cost while leaving accidental gains in the numerator.

A profitable business whose founder teaches 60 unpaid hours a month may be a respectable owner-operated livelihood; it is not necessarily an acquirable organization at the advertised margin. A loss-making firm expensing an unusually large 1-time product build may have a credible path to profitability-but only if future maintenance and growth needs are modeled honestly. \textbf{Normalize reality, not hope.}

\subsection{37. Cohort and vintage analysis reveals hidden economics}


Aggregate financial statements can hide deterioration across customer generations. Study customer cohorts by acquisition quarter, onboarding period, region and pricing plan. In lending, a vintage describes exposures originated in a particular period; in subscriptions it often means a signup cohort; in education it may be a program intake. Compare retention, support cost, refunds, collection, delivery quality and contribution on consistent definitions.

A software startup's revenue may grow because it adds customers faster than older customers churn. A school may report an excellent completion rate among learners still active, quietly excluding withdrawals. A financing company may report rising income before the losses on recent loans have matured. Cohort analysis shifts the question from \textbf{how large are we today?} to \textbf{does each new generation of business carry economics at least as good as the last?}

\subsection{38. Accrual quality and the cash-conversion check}


Over several periods, compare net income and operating cash flow after normalizing growth in working capital. Persistent large accounting profits with weak collection may have an innocuous explanation-rapid credit growth-or indicate aggressive recognition and poor credit quality. Track \textbf{DSO, disputed invoices, credit notes after period-end, customer prepayments, and changes in contract assets}. Understand seasonality and sales concentration before alleging misconduct.

Cash conversion is informative but not sufficient: a company can generate strong operating cash by delaying suppliers unsustainably, taking customer advances for undelivered services, selling receivables, or cutting maintenance. Conversely, responsible investment in working capital can temporarily depress operating cash during profitable expansion. Forensics requires explaining \textit{sources and sustainability}, not merely calculating a percentage.

\subsection{39. The 3-statement founder examination}


A financially literate founder should be able to answer without searching a dashboard: Why did profit differ from cash? Which assets generate returns? Which liabilities have hard settlement dates? How much capital is tied in operating working capital? Which customer segments generate contribution? What cash must be preserved to finish sold work? What capital expenditures are essential rather than optional? Which commitments would survive a demand shock? Which reported adjustments depend on subjective management claims?

If management cannot produce these answers, a valuation or funding discussion is premature. The company may still be good; its economic reality has not yet been represented well enough to finance responsibly.

